% STATUS: written 2026-11-05 (Hyers 1941 and Miura-Hirasawa-Takahasi 2006
% read in full; excerpts in excerpts/Hyers1941Stability.txt and
% excerpts/MiuraHirasawaTakahasi2006.txt).
% Additions of 2026-11-07 from Roh--Jung--Chang 2009 (pp. 1975--1986) and
% Konig--Milman 2014 (pp. 405--407).
% Sources: Hyers1941Stability, MiuraHirasawaTakahasi2006,
% RohJungChang2009Triple, KonigMilman2014Derivation.
\section{Approximate derivations and the stability of functional equations}\label{sec:banach-operator:stability}

Section~\ref{sec:banach-operator:commutator-ranges} studied exact
derivations. This section asks whether a map that satisfies the equations
only approximately is near a map that satisfies them exactly
\cite[p.~222]{Hyers1941Stability}. Ulam posed the question at the
University of Wisconsin \cite[p.~222]{Hyers1941Stability}. Let $E$ and
$E'$ be Banach spaces and $\delta$ a positive number; a transformation
$f$ of $E$ into $E'$ with $\|f(x+y)-f(x)-f(y)\|<\delta$ for all $x,y$ in
$E$ is called \emph{$\delta$-linear}. Must it stay within a fixed
distance of a linear one \cite[p.~222]{Hyers1941Stability}? Hyers proved
that one always exists in 1941 \cite[Theorem~1, p.~223]{Hyers1941Stability}.

\begin{theorem}[{\cite[Theorem~1, p.~223]{Hyers1941Stability}}]\label{thm:banach:hyers}
Let $E$ and $E'$ be Banach spaces and let $f(x)$ be a $\delta$-linear
transformation of $E$ into $E'$. Then the limit
$l(x)=\lim_{n\to\infty}f(2^nx)/2^n$ exists for each $x$ in $E$, $l(x)$ is
a linear transformation, and the inequality $\|f(x)-l(x)\|\le\delta$ is
true for all $x$ in $E$. Moreover $l(x)$ is the only linear transformation
satisfying this inequality.
\end{theorem}

Hyers added a continuity observation
\cite[Theorem~2, p.~224]{Hyers1941Stability}.

\begin{theorem}[{\cite[Theorem~2, p.~224]{Hyers1941Stability}}]\label{thm:banach:hyerscont}
If under the hypotheses of Theorem~\ref{thm:banach:hyers} we suppose that
$f(x)$ is continuous at a single point $y$ of $E$, then $l(x)$ is
continuous everywhere in $E$.
\end{theorem}

Theorem~\ref{thm:banach:hyers} is called the Hyers--Ulam stability of the
additive Cauchy equation \cite[p.~524]{MiuraHirasawaTakahasi2006}. Later
writers widened the control function, and Miura, Hirasawa and Takahasi
brought the idea to ring derivations
\cite[p.~524]{MiuraHirasawaTakahasi2006}.

\begin{definition}[{\cite[p.~524]{MiuraHirasawaTakahasi2006}}]\label{def:banach:approxder}
Let $A$ be a Banach algebra, $p\ge 0$ and $\varepsilon\ge 0$. A mapping
$f:A\to A$ is an \emph{approximate ring derivation} if
$\|f(a+b)-f(a)-f(b)\|\le\varepsilon(\|a\|^p+\|b\|^p)$ for all $a,b$ in
$A$, and $\|f(ab)-af(b)-f(a)b\|\le\varepsilon\|a\|^p\|b\|^p$ for all
$a,b$ in $A$.
\end{definition}

The history of such control functions runs through Rassias, with error
$\varepsilon(\|x\|^p+\|y\|^p)$ for $0\le p<1$, and Gajda, who settled
$1<p$ and gave an example that stability fails at $p=1$; Bourgin proved
the superstability of ring homomorphisms, where the map is exactly a
homomorphism \cite[pp.~524--526]{MiuraHirasawaTakahasi2006}. For
derivations on general algebras, an approximate derivation need not be an
exact one. On $C(X)$ over a compact Hausdorff space $X$ without isolated
points, the map $f(a)(x)=a(x)\log|a(x)|$ satisfies the derivation identity
\cite[pp.~523--524]{MiuraHirasawaTakahasi2006}. It is a slight
modification of an example of Semrl whose idea is due to
Johnson, and it differs from every additive map by an unbounded term
\cite[pp.~523--524]{MiuraHirasawaTakahasi2006}. Superstability holds once
the algebra has no zero divisor from either side
\cite[Theorem~1, p.~526]{MiuraHirasawaTakahasi2006}. A Banach algebra is
\emph{without order} if no nonzero $x_0$ and $y_0$ satisfy
$x_0A=Ay_0=\{0\}$; semisimple commutative Banach algebras and Banach
algebras with approximate identities are without order
\cite[p.~526]{MiuraHirasawaTakahasi2006}.

\begin{theorem}[{\cite[Theorem~1, p.~526]{MiuraHirasawaTakahasi2006}}]\label{thm:banach:superstable}
Suppose $A$ is a Banach algebra without order. If $f:A\to A$ is a mapping
satisfying the two inequalities of
Definition~\ref{def:banach:approxder} for some $\varepsilon\ge 0$ and
$p\ge 0$, $p\ne 1$, then $f$ is a ring derivation.
\end{theorem}

Theorem~\ref{thm:banach:superstable} answers a question raised by
Watanabe \cite[p.~526]{MiuraHirasawaTakahasi2006}. On functions it has a zero
consequence \cite[Corollary~3, p.~526]{MiuraHirasawaTakahasi2006}.

\begin{corollary}[{\cite[Corollary~3, p.~526]{MiuraHirasawaTakahasi2006}}]\label{cor:banach:superzero}
Suppose $A$ is a semisimple commutative Banach algebra with the maximal
ideal space $M_A$ without isolated points. If $f:A\to A$ satisfies the two
inequalities of Definition~\ref{def:banach:approxder} for some
$\varepsilon\ge 0$, $p\ge 0$, $p\ne 1$, then $f$ is identically zero.
\end{corollary}

Corollary~\ref{cor:banach:superzero} is the approximate form of the
Singer--Wermer theorem of
Section~\ref{sec:banach-operator:singer-wermer}: where the exact theory
gives only the zero derivation, the approximate theory gives only the zero
map. The chapter closed as it opened, on commutative Banach algebras
\cite[Corollary~3, p.~526]{MiuraHirasawaTakahasi2006}.

The same question was carried to higher derivations and to the Jordan triple
product \cite[p.~1975]{RohJungChang2009Triple}.
Roh, Jung and Chang control a sequence of maps by an inequality for the
Jordan triple product $xyx$
\cite[p.~1976]{RohJungChang2009Triple}. The rule reads
$f_n(xyx)=\sum_{i+j+k=n}f_i(x)f_j(y)f_k(x)$.

\begin{theorem}[{\cite[Theorem~2.6, p.~1984]{RohJungChang2009Triple}}]\label{thm:banach:roh-triple-radical}
Let $A$ be a commutative semiprime Banach algebra. Suppose that a sequence
$F=\{f_0,f_1,\dots,f_m\}$ of continuous operators on $A$, with $f_0$ the
identity, satisfies the inequalities of the paper for the Jordan triple rule.
Then $F$ maps $A$ into its radical.
\end{theorem}

On a commutative semisimple algebra the same hypotheses force the map to be
zero \cite[Corollary~2.8, p.~1986]{RohJungChang2009Triple}. This is the
higher-order form of Corollary~\ref{cor:banach:superzero}
\cite[p.~1986]{RohJungChang2009Triple}.

Continuity also enters the stability question directly
\cite[p.~405]{KonigMilman2014Derivation}. K\"onig and Milman
ask which operators on $C^k(\mathbb{R})$ satisfy the Leibniz rule up to a
small error term, and the answer is a rigidity theorem
\cite[p.~405]{KonigMilman2014Derivation}. The error term is a function $B$ of
the derivatives of order below $k$ only
\cite[p.~407]{KonigMilman2014Derivation}.

\begin{theorem}[{\cite[Theorem~1, p.~407]{KonigMilman2014Derivation}}]\label{thm:banach:konig-milman}
Let $k$ be a positive integer, $T\colon C^k(\mathbb{R})\to C(\mathbb{R})$ an
operator, and $B$ a function such that
$T(fg)=Tf\cdot g+f\cdot Tg+B(x,f,\dots,f^{(k-1)},g,\dots,g^{(k-1)})$ for all
$f,g$ and all $x$. Suppose that $T$ annihilates every polynomial of degree at
most $k-1$. Then $T$ is a multiple of the $k$-th derivative, $Tf=d\,f^{(k)}$
for a continuous function $d$, and $B$ is the corresponding sum from the
$k$-th order Leibniz rule.
\end{theorem}

The theorem is a stability result for the $k$-th derivative
\cite[p.~407]{KonigMilman2014Derivation}. For $k=2$ on $\mathbb{R}^n$ the same authors describe the Laplacian by a
similar equation \cite[p.~405]{KonigMilman2014Derivation}. Two further requirements enter there: the operator is unchanged under maps
that preserve the inner product, and it kills the polynomials of lower order
\cite[p.~405]{KonigMilman2014Derivation}.

Two later works measure how far a derivation can reach. Kumar and Wirth use
derivations as the tool for coarse rigidity of von Neumann algebras
\cite[Abstract, p.~1]{KumarWirth2026}. Lupini compares the length a derivation
needs with the length an automorphism needs on a unital $C^{*}$-algebra
\cite[Abstract, p.~1]{Lupini2026Length}. Park, in the same setting, treats
permuting triderivations together with the permuting trihomomorphisms they
extend \cite[Abstract, p.~1]{Park2020Trideriv}.

The stability and superstability of generalized $(\theta,\phi)$-derivations on non-Archimedean algebras are proved by a fixed-point argument on the additive Cauchy functional equation~\cite{WlStabilityAndSuperstability2019}.

An approximate derivation of a Banach algebra whose defect lies in the radical is shown to differ from a derivation by an element of the radical, which is the radical-range form of the stability theorem~\cite{WlApproximateDerivationsWithTheR2015}.

The Hyers-Ulam-Rassias stability of a functional equation on Banach spaces is proved, and the stability of homomorphisms and derivations on $C$-ternary algebras follows from it~\cite{WlStabilityOfHomomorphismsAndDer2019}.

Under stated conditions a double Jordan derivation is approximated by a function near it in a $p$-Banach algebra, which is the generalized-stability form of the Jordan-derivation result~\cite{WlFunctionsNearSomeDoubleJordanD2019}.
